\documentclass[]{article}
\usepackage[utf8]{inputenc}
\usepackage{amsmath, amssymb, amsthm}
\usepackage{geometry}
\usepackage{xcolor}
\usepackage[most]{tcolorbox}
\usepackage{enumitem}
\usepackage{fancyhdr}

% Page setup
\geometry{margin=1in}
\pagestyle{fancy}
\fancyhf{}
\rhead{AMC12A 2016 Problem 15 Analysis}
\lhead{\today}

% Color definitions
\definecolor{problemconfig}{RGB}{230, 242, 255}
\definecolor{strategic}{RGB}{255, 230, 242}
\definecolor{tactical}{RGB}{242, 255, 230}
\definecolor{techniques}{RGB}{255, 242, 230}
\definecolor{verification}{RGB}{255, 230, 230}
\definecolor{solution}{RGB}{230, 255, 255}
\definecolor{competition}{RGB}{242, 242, 255}
\definecolor{gold}{RGB}{218, 165, 32}

% Box styles
\newtcolorbox{problemconfigbox}{
    colback=problemconfig,
    colframe=blue,
    title=Problem Configuration Analysis,
    fonttitle=\bfseries,
    arc=3pt,
    breakable,
    boxrule=1pt
}

\newtcolorbox{strategicbox}{
    colback=strategic,
    colframe=purple,
    title=Strategic Framework Application,
    fonttitle=\bfseries,
    arc=3pt,
    breakable,
    boxrule=1pt
}

\newtcolorbox{tacticalbox}{
    colback=tactical,
    colframe=green,
    title=Tactical Methodology Selection,
    fonttitle=\bfseries,
    arc=3pt,
    breakable,
    boxrule=1pt
}

\newtcolorbox{techniquesbox}{
    colback=techniques,
    colframe=orange,
    title=Conversion Techniques Application,
    fonttitle=\bfseries,
    arc=3pt,
    breakable,
    boxrule=1pt
}

\newtcolorbox{verificationbox}{
    colback=verification,
    colframe=red,
    title=Verification Strategy Design,
    fonttitle=\bfseries,
    arc=3pt,
    breakable,
    boxrule=1pt
}

\newtcolorbox{solutionbox}{
    colback=solution,
    colframe=cyan,
    title=Solution Roadmap,
    fonttitle=\bfseries,
    arc=3pt,
    breakable,
    boxrule=1pt
}

\newtcolorbox{competitionbox}{
    colback=competition,
    colframe=purple,
    title=Competition Strategy Integration,
    fonttitle=\bfseries,
    arc=3pt,
    breakable,
    boxrule=1pt
}

\newtcolorbox{keyinsightbox}{
    colback=yellow!20,
    colframe=gold,
    title=Key Insight,
    fonttitle=\bfseries,
    arc=3pt,
    boxrule=2pt,
    leftrule=5pt,
    breakable
}

\newtcolorbox{warningbox}{
    colback=red!10,
    colframe=red!50,
    title=Warning: Common Pitfall,
    fonttitle=\bfseries,
    arc=3pt,
    boxrule=2pt,
    leftrule=5pt,
    breakable
}

\newtcolorbox{tipbox}{
    colback=green!10,
    colframe=green!50,
    title=Pro Tip,
    fonttitle=\bfseries,
    arc=3pt,
    boxrule=2pt,
    leftrule=5pt,
    breakable
}

\begin{document}

\title{AMC12/COMC Strategic Problem Analysis\\2016 AMC 12A Problem 15}
\author{Competition Math Framework}
\date{\today}
\maketitle

\section*{Problem Statement}

Circles with centers $P, Q$ and $R$, having radii $1, 2$ and $3$, respectively, lie on the same side of line $l$ and are tangent to $l$ at $P', Q'$ and $R'$, respectively, with $Q'$ between $P'$ and $R'$. The circle with center $Q$ is externally tangent to each of the other two circles. What is the area of triangle $PQR$?

\[\textbf{(A) } 0\qquad \textbf{(B) } \sqrt{\frac{2}{3}}\qquad\textbf{(C) } 1\qquad\textbf{(D) } \sqrt{6}-\sqrt{2}\qquad\textbf{(E) }\sqrt{\frac{3}{2}}\]

\begin{problemconfigbox}
\subsection*{A. Problem Classification}
\begin{itemize}[leftmargin=*]
    \item \textbf{Problem Type}: To Find - calculate the area of triangle $PQR$
    \item \textbf{Difficulty Band}: Mid-band (Problem 15/25, transitioning to late-band)
    \item \textbf{Competition Context}: AMC12 (75 min, 25 problems) - Target time 4-5 minutes
    \item \textbf{Mathematical Domain}: Plane Geometry with strong algebraic computation component
    \item \textbf{Source Context}: 2016 AMC 12A Problem 15 - historically correct answer rate ~15-20\%
\end{itemize}

\subsection*{B. Problem Structure Identification}
\begin{itemize}[leftmargin=*]
    \item \textbf{Given Information}: 
    \begin{itemize}
        \item Three circles with radii 1, 2, 3 centered at $P$, $Q$, $R$ respectively
        \item All circles tangent to line $l$ from same side
        \item Tangent points $P'$, $Q'$, $R'$ with $Q'$ between $P'$ and $R'$
        \item Circle $Q$ externally tangent to both circles $P$ and $R$
    \end{itemize}
    \item \textbf{Constraints}: 
    \begin{itemize}
        \item External tangency implies $PQ = 1 + 2 = 3$ and $QR = 2 + 3 = 5$
        \item Tangency to line implies heights are the radii (1, 2, 3)
        \item Ordering constraint: $Q'$ between $P'$ and $R'$
    \end{itemize}
    \item \textbf{Objective}: Find area of triangle $PQR$ from answer choices
    \item \textbf{Hidden Assumptions}: 
    \begin{itemize}
        \item Configuration is uniquely determined by constraints
        \item Triangle exists (non-collinear points)
        \item Answer involves nested radicals (based on choices)
    \end{itemize}
    \item \textbf{Answer Choice Leverage}: 
    \begin{itemize}
        \item (A) is trivial case - immediately eliminable (tangency conditions imply non-degeneracy)
        \item (B), (C), (E) are relatively simple expressions
        \item (D) $\sqrt{6} - \sqrt{2}$ is most complex - likely answer if problem is difficult
        \item Can verify by checking if calculated area matches a choice
    \end{itemize}
\end{itemize}

\subsection*{C. Strategic Orientation}
\begin{itemize}[leftmargin=*]
    \item \textbf{Hypothesis}: The tangency conditions and external tangency determine unique configuration
    \item \textbf{Conclusion}: Area can be computed via coordinate geometry or trapezoid decomposition
    \item \textbf{Notation}: 
    \begin{itemize}
        \item Centers: $P$, $Q$, $R$
        \item Radii: $r_P = 1$, $r_Q = 2$, $r_R = 3$
        \item Tangent points: $P'$, $Q'$, $R'$ on line $l$
        \item Distances: $PQ = 3$, $QR = 5$, $PR = ?$
    \end{itemize}
    \item \textbf{Intuitive Approach}: Use trapezoid decomposition method - break into manageable pieces
    \item \textbf{Cross-Domain Potential}: 
    \begin{itemize}
        \item Coordinate geometry: Place line $l$ as $x$-axis
        \item Pythagorean theorem: Find horizontal distances between tangent points
        \item Heron's formula: After finding all three side lengths
    \end{itemize}
\end{itemize}
\end{problemconfigbox}

\begin{strategicbox}
\subsection*{A. Psychological Strategies}
\begin{itemize}[leftmargin=*]
    \item \textbf{Mental Toughness}: This is a mid-band problem requiring careful calculation. Multiple approaches exist - if one gets messy, switch methods.
    \item \textbf{Creativity Cultivation}: The trapezoid decomposition approach is elegant - visualize breaking the configuration into trapezoids and triangles.
    \item \textbf{Iterative Refinement}: If initial calculation produces messy algebra, simplify by choosing better coordinate system or decomposition strategy.
\end{itemize}

\subsection*{B. Investigation Strategies}
\begin{itemize}[leftmargin=*]
    \item \textbf{Startup Orientation}: Draw accurate diagram showing three circles tangent to line with proper ordering
    \item \textbf{Get Your Hands Dirty}: Calculate $P'Q'$ and $Q'R'$ using Pythagorean theorem on trapezoids
    \item \textbf{Penultimate Step}: After finding all trapezoid areas, final step is simple addition/subtraction
    \item \textbf{Wishful Thinking}: "If only we knew $PR$..." - but we can find it via Pythagorean theorem
    \item \textbf{Make It Easier}: Instead of finding $PR$ directly, use trapezoid area method to avoid that calculation
    \item \textbf{Conjecture via Small Cases}: Not applicable - unique configuration
    \item \textbf{Visualization}: Critical - draw perpendiculars from centers to line to form trapezoids
\end{itemize}

\subsection*{C. Late-Band Strategic Playbook}
\begin{itemize}[leftmargin=*]
    \item \textbf{Answer-Choice Leverage}: Answer (D) $\sqrt{6} - \sqrt{2}$ suggests problem involves nested radicals - consistent with Pythagorean calculations yielding $\sqrt{8}$ and $\sqrt{24}$
    \item \textbf{Make It Easier}: Not used - problem structure is already clean
    \item \textbf{Penultimate Step}: Final calculation is adding trapezoid areas and subtracting large trapezoid
    \item \textbf{Wishful Thinking}: Not needed - direct approach works
    \item \textbf{Conjecture via Small Cases}: Verify answer makes sense: area should be positive and reasonable given side lengths 3, 5
\end{itemize}

\subsection*{D. Argument Construction Strategy}
\begin{itemize}[leftmargin=*]
    \item \textbf{Direct Proof}: Yes - calculate areas using trapezoid decomposition method
    \item \textbf{Proof by Contradiction}: Not applicable
    \item \textbf{Mathematical Induction}: Not applicable
    \item \textbf{Stylistic Conventions}: WLOG can assume coordinate system with $l$ as $x$-axis
\end{itemize}

\subsection*{E. Cross-Domain Translation}
\begin{itemize}[leftmargin=*]
    \item \textbf{Draw a Picture}: Essential - diagram reveals trapezoid structure
    \item \textbf{Recast the Problem}: "Find area" becomes "decompose into trapezoids and calculate"
    \item \textbf{Change Your Point of View}: Coordinate geometry approach: Shoelace formula after finding coordinates
\end{itemize}
\end{strategicbox}

\begin{tacticalbox}
\subsection*{A. Core Mathematical Tactics}
\begin{itemize}[leftmargin=*]
    \item \textbf{Symmetry}: Configuration has no obvious symmetry, but trapezoid method exploits structural regularity
    \item \textbf{The Extreme Principle}: Not applicable
    \item \textbf{The Pigeonhole Principle}: Not applicable
    \item \textbf{Invariants}: The heights (radii) remain constant throughout calculation
\end{itemize}

\subsection*{B. Crossover Tactics}
\begin{itemize}[leftmargin=*]
    \item \textbf{Graph Theory}: Not applicable
    \item \textbf{Complex Numbers}: Not applicable (though could represent as complex plane)
    \item \textbf{Generating Functions}: Not applicable
    \item \textbf{Coordinate Geometry}: Highly applicable - primary alternative approach using Shoelace theorem
\end{itemize}
\end{tacticalbox}

\begin{techniquesbox}
\subsection*{Recognition Index Analysis}
\textbf{Primary Recognition Triggers (1-3 seconds):}
\begin{enumerate}
    \item \textbf{Circles tangent to line} $\rightarrow$ Heights equal radii (perpendicular from center to tangent point)
    \item \textbf{External tangency} $\rightarrow$ Distance between centers equals sum of radii
    \item \textbf{Find area of triangle} $\rightarrow$ Multiple approaches: coordinate geometry, Heron's, decomposition
    \item \textbf{Answer choices with nested radicals} $\rightarrow$ Expect Pythagorean theorem calculations
\end{enumerate}

\subsection*{A. High-Probability Techniques Selected}
\begin{enumerate}[leftmargin=*]
    \item \textbf{Coordinate Geometry Conversion}: 
    \begin{itemize}
        \item Place $P'$ at origin, line $l$ as $x$-axis
        \item Centers: $P = (0, 1)$, $Q = (P'Q', 2)$, $R = (P'Q' + Q'R', 3)$
        \item Find $P'Q'$ and $Q'R'$ using Pythagorean theorem
        \item Apply Shoelace formula for area
    \end{itemize}
    
    \item \textbf{Pythagorean Theorem Application}:
    \begin{itemize}
        \item In trapezoid $PP'Q'Q$: $P'Q'^2 + (r_Q - r_P)^2 = PQ^2$
        \item $P'Q'^2 + 1^2 = 3^2 \Rightarrow P'Q' = \sqrt{8} = 2\sqrt{2}$
        \item In trapezoid $QQ'R'R$: $Q'R'^2 + (r_R - r_Q)^2 = QR^2$
        \item $Q'R'^2 + 1^2 = 5^2 \Rightarrow Q'R' = \sqrt{24} = 2\sqrt{6}$
    \end{itemize}
    
    \item \textbf{Trapezoid Area Decomposition}:
    \begin{itemize}
        \item Area of trapezoid with parallel sides $a$, $b$ and height $h$: $\frac{(a+b)h}{2}$
        \item $[PP'Q'Q] = \frac{(r_P + r_Q) \cdot P'Q'}{2} = \frac{(1+2) \cdot 2\sqrt{2}}{2} = 3\sqrt{2}$
        \item $[QQ'R'R] = \frac{(r_Q + r_R) \cdot Q'R'}{2} = \frac{(2+3) \cdot 2\sqrt{6}}{2} = 5\sqrt{6}$
        \item $[PP'R'R] = \frac{(r_P + r_R) \cdot P'R'}{2} = \frac{(1+3)(2\sqrt{2}+2\sqrt{6})}{2} = 4\sqrt{2} + 4\sqrt{6}$
    \end{itemize}
    
    \item \textbf{Area Addition/Subtraction}:
    \begin{itemize}
        \item Key equation: $[PP'Q'Q] + [QQ'R'R] = [PQR] + [PP'R'R]$
        \item Therefore: $[PQR] = [PP'Q'Q] + [QQ'R'R] - [PP'R'R]$
        \item $[PQR] = 3\sqrt{2} + 5\sqrt{6} - (4\sqrt{2} + 4\sqrt{6})$
        \item $[PQR] = 3\sqrt{2} + 5\sqrt{6} - 4\sqrt{2} - 4\sqrt{6}$
        \item $[PQR] = -\sqrt{2} + \sqrt{6} = \sqrt{6} - \sqrt{2}$
    \end{itemize}
\end{enumerate}

\subsection*{B. Medium-Probability Techniques}
\begin{enumerate}[leftmargin=*]
    \item \textbf{Heron's Formula}: 
    \begin{itemize}
        \item Would require finding $PR$ first via Pythagorean theorem
        \item $PR^2 = (P'R')^2 + (r_R - r_P)^2 = (2\sqrt{2}+2\sqrt{6})^2 + 4$
        \item Leads to nested radical: $PR = 2\sqrt{9 + 4\sqrt{3}}$
        \item More complex but verifiable approach
    \end{itemize}
\end{enumerate}

\subsection*{C. Technique Selection Rationale}
\begin{itemize}[leftmargin=*]
    \item \textbf{Primary Method}: Trapezoid decomposition - clean, direct, avoids finding $PR$
    \item \textbf{Alternative Method}: Coordinate geometry with Shoelace - equally clean
    \item \textbf{Backup Method}: Heron's formula - more computation but verifies answer
    \item \textbf{Why Chosen}: Trapezoid method is fastest and least error-prone for competition
\end{itemize}

\subsection*{D. Integration Strategy}
\begin{itemize}[leftmargin=*]
    \item All three methods should yield $\sqrt{6} - \sqrt{2}$
    \item Trapezoid method integrates: Pythagorean theorem + area formula + algebraic manipulation
    \item No conflicts between techniques - they complement each other
\end{itemize}
\end{techniquesbox}

\begin{verificationbox}
\subsection*{A. Verification Level Selection}
Based on decision tree for mid-band problem with computational complexity:

\textbf{Selected Level: Level 4 - Standard Verification}
\begin{itemize}[leftmargin=*]
    \item Difficulty: Mid-band (P15)
    \item Computation: Moderate (multiple Pythagorean calculations)
    \item Confidence: Should be high after clean trapezoid calculation
    \item Time available: 1-1.5 minutes for verification
\end{itemize}

\subsection*{B. Specific Verification Methods}
\begin{enumerate}[leftmargin=*]
    \item \textbf{Dimensional Analysis}:
    \begin{itemize}
        \item Area should have dimension [length]$^2$
        \item $\sqrt{6} - \sqrt{2}$ is dimensionless number, multiply by appropriate unit squared
        \item Check: Expression is positive ($\sqrt{6} \approx 2.45 > \sqrt{2} \approx 1.41$) $\checkmark$
    \end{itemize}
    
    \item \textbf{Reasonableness Check}:
    \begin{itemize}
        \item Triangle has sides $PQ = 3$, $QR = 5$
        \item Rough upper bound: $\frac{1}{2} \cdot 3 \cdot 5 = 7.5$
        \item Computed: $\sqrt{6} - \sqrt{2} \approx 2.45 - 1.41 = 1.04$
        \item Reasonable - triangle is not very "tall" $\checkmark$
    \end{itemize}
    
    \item \textbf{Answer Choice Matching}:
    \begin{itemize}
        \item Computed answer exactly matches choice (D)
        \item No other choice is close numerically $\checkmark$
    \end{itemize}
    
    \item \textbf{Algebraic Verification}:
    \begin{itemize}
        \item Re-check: $3\sqrt{2} + 5\sqrt{6} - 4\sqrt{2} - 4\sqrt{6}$
        \item Group terms: $(3-4)\sqrt{2} + (5-4)\sqrt{6} = -\sqrt{2} + \sqrt{6}$ $\checkmark$
    \end{itemize}
\end{enumerate}

\subsection*{C. Specialized Verification}
\textbf{Geometry-Specific Checks:}
\begin{itemize}[leftmargin=*]
    \item Verify $P'Q' = 2\sqrt{2}$: $\sqrt{3^2 - 1^2} = \sqrt{8} = 2\sqrt{2}$ $\checkmark$
    \item Verify $Q'R' = 2\sqrt{6}$: $\sqrt{5^2 - 1^2} = \sqrt{24} = 2\sqrt{6}$ $\checkmark$
    \item Triangle inequality: $3 + 5 > PR$, $|5-3| < PR$ (should be satisfied)
\end{itemize}

\subsection*{D. Confidence Calibration}
\begin{itemize}[leftmargin=*]
    \item \textbf{Initial Confidence}: 85\% (clean calculation, matches answer choice)
    \item \textbf{After Verification}: 95\% (all checks pass, answer is exact match)
    \item \textbf{Flag for Review}: No - high confidence, move to next problem
\end{itemize}
\end{verificationbox}

\begin{solutionbox}
\subsection*{Complete Solution Roadmap}

\textbf{Time Estimate: 4-5 minutes}

\subsubsection*{Step 1: Setup and Diagram (30 seconds)}
\begin{itemize}[leftmargin=*]
    \item Draw line $l$ horizontally
    \item Draw three circles tangent to $l$ from above
    \item Mark radii: 1, 2, 3 for circles P, Q, R
    \item Note ordering: $P'$, $Q'$, $R'$ left to right
    \item Mark tangency: Q externally tangent to P and R
\end{itemize}

\textbf{Checkpoint}: Diagram should clearly show configuration, verify $PQ = 3$, $QR = 5$

\subsubsection*{Step 2: Find Horizontal Distances (1.5 minutes)}
\begin{itemize}[leftmargin=*]
    \item For trapezoid $PP'Q'Q$:
    \begin{itemize}
        \item Vertical difference: $|r_Q - r_P| = |2-1| = 1$
        \item Hypotenuse: $PQ = 3$
        \item Horizontal: $P'Q' = \sqrt{3^2 - 1^2} = \sqrt{8} = 2\sqrt{2}$
    \end{itemize}
    
    \item For trapezoid $QQ'R'R$:
    \begin{itemize}
        \item Vertical difference: $|r_R - r_Q| = |3-2| = 1$
        \item Hypotenuse: $QR = 5$
        \item Horizontal: $Q'R' = \sqrt{5^2 - 1^2} = \sqrt{24} = 2\sqrt{6}$
    \end{itemize}
\end{itemize}

\textbf{Checkpoint}: $P'Q' = 2\sqrt{2}$, $Q'R' = 2\sqrt{6}$, therefore $P'R' = 2\sqrt{2} + 2\sqrt{6}$

\subsubsection*{Step 3: Calculate Trapezoid Areas (1.5 minutes)}
\begin{itemize}[leftmargin=*]
    \item Area of $PP'Q'Q$:
    \[[\text{PP'Q'Q}] = \frac{(r_P + r_Q) \cdot P'Q'}{2} = \frac{(1+2)(2\sqrt{2})}{2} = 3\sqrt{2}\]
    
    \item Area of $QQ'R'R$:
    \[[\text{QQ'R'R}] = \frac{(r_Q + r_R) \cdot Q'R'}{2} = \frac{(2+3)(2\sqrt{6})}{2} = 5\sqrt{6}\]
    
    \item Area of $PP'R'R$:
    \begin{align*}
    [\text{PP'R'R}] &= \frac{(r_P + r_R) \cdot P'R'}{2}\\
    &= \frac{(1+3)(2\sqrt{2}+2\sqrt{6})}{2}\\
    &= \frac{4(2\sqrt{2}+2\sqrt{6})}{2}\\
    &= 2(2\sqrt{2}+2\sqrt{6})\\
    &= 4\sqrt{2} + 4\sqrt{6}
    \end{align*}
\end{itemize}

\textbf{Checkpoint}: Three trapezoid areas calculated, ready for final combination

\subsubsection*{Step 4: Apply Area Equation (1 minute)}
Key relationship: $[PP'Q'Q] + [QQ'R'R] = [PQR] + [PP'R'R]$

Therefore:
\begin{align*}
[PQR] &= [PP'Q'Q] + [QQ'R'R] - [PP'R'R]\\
&= 3\sqrt{2} + 5\sqrt{6} - (4\sqrt{2} + 4\sqrt{6})\\
&= 3\sqrt{2} + 5\sqrt{6} - 4\sqrt{2} - 4\sqrt{6}\\
&= (3-4)\sqrt{2} + (5-4)\sqrt{6}\\
&= -\sqrt{2} + \sqrt{6}\\
&= \sqrt{6} - \sqrt{2}
\end{align*}

\textbf{Checkpoint}: Answer matches choice (D)

\subsubsection*{Step 5: Verification (30 seconds)}
\begin{itemize}[leftmargin=*]
    \item Check sign: $\sqrt{6} > \sqrt{2}$, so area is positive $\checkmark$
    \item Numerical check: $\sqrt{6} - \sqrt{2} \approx 2.45 - 1.41 = 1.04$ $\checkmark$
    \item Answer choice match: (D) $\checkmark$
\end{itemize}

\subsection*{Alternative Approaches}
\begin{enumerate}[leftmargin=*]
    \item \textbf{Coordinate Geometry + Shoelace}:
    \begin{itemize}
        \item Place $P = (0,1)$, $Q = (2\sqrt{2}, 2)$, $R = (2\sqrt{2}+2\sqrt{6}, 3)$
        \item Apply Shoelace formula
        \item Time: Similar to trapezoid method
    \end{itemize}
    
    \item \textbf{Heron's Formula}:
    \begin{itemize}
        \item Find $PR = 2\sqrt{9+4\sqrt{3}}$
        \item Apply Heron's with sides 3, 5, $2\sqrt{9+4\sqrt{3}}$
        \item Time: 6-7 minutes (more complex)
    \end{itemize}
\end{enumerate}

\subsection*{Pitfall Guide}
\begin{enumerate}[leftmargin=*]
    \item \textbf{Arithmetic Error}: Miscalculating $\sqrt{8}$ or $\sqrt{24}$
    \begin{itemize}
        \item Prevention: Double-check simplification of radicals
        \item Recovery: Verify with calculator if allowed
    \end{itemize}
    
    \item \textbf{Area Equation Error}: Wrong combination of trapezoid areas
    \begin{itemize}
        \item Prevention: Draw diagram clearly showing which areas add/subtract
        \item Recovery: Re-derive from first principles
    \end{itemize}
    
    \item \textbf{Sign Error}: Getting $\sqrt{2} - \sqrt{6}$ (negative area)
    \begin{itemize}
        \item Prevention: Always check final answer is positive
        \item Recovery: Review area combination equation
    \end{itemize}
    
    \item \textbf{Pythagorean Theorem Application}: Using wrong sides
    \begin{itemize}
        \item Prevention: Label trapezoid clearly with known quantities
        \item Recovery: Verify vertical difference is radius difference
    \end{itemize}
\end{enumerate}

\subsection*{Time Management Checkpoints}
\begin{itemize}[leftmargin=*]
    \item \textbf{1 minute}: If still setting up diagram, speed up or skip
    \item \textbf{3 minutes}: Should have horizontal distances calculated
    \item \textbf{5 minutes}: If not near final answer, consider flagging and moving on
    \item \textbf{Pivot Indicator}: If algebra becomes messy with terms $>3$ radicals, switch to coordinate approach
\end{itemize}
\end{solutionbox}

\begin{competitionbox}
\subsection*{A. Time Management Strategy}
\begin{itemize}[leftmargin=*]
    \item \textbf{Target Time}: 4-5 minutes (appropriate for P15)
    \item \textbf{Maximum Time}: 6 minutes (includes verification)
    \item \textbf{Actual Breakdown}:
    \begin{itemize}
        \item Reading + diagram: 30 seconds
        \item Finding distances: 1.5 minutes
        \item Computing areas: 1.5 minutes
        \item Final calculation: 1 minute
        \item Verification: 30 seconds
    \end{itemize}
\end{itemize}

\subsection*{B. Problem Prioritization}
\begin{itemize}[leftmargin=*]
    \item \textbf{Accessibility}: HIGH - clear geometric setup, standard techniques
    \item \textbf{Difficulty}: MEDIUM-HIGH (P15) - requires careful calculation
    \item \textbf{Point Value}: 6 points (standard AMC12)
    \item \textbf{Strategic Position}: Should attempt after securing P1-14
    \item \textbf{Domain Comfort}: If strong in geometry, prioritize; if weak, consider skipping
\end{itemize}

\subsection*{C. Risk Management}
\begin{itemize}[leftmargin=*]
    \item \textbf{Confidence Threshold}: 90\%+ after verification - submit and move on
    \item \textbf{Medium Confidence (70-90\%)}: Flag for review if time permits
    \item \textbf{Low Confidence (<70\%)}: Make educated guess between (C) and (D)
\end{itemize}

\subsection*{D. Abandonment Criteria}
\textbf{Soft Abandon Triggers (Consider switching):}
\begin{itemize}[leftmargin=*]
    \item After 5 minutes without clear path to answer
    \item Computational errors accumulating (2+ mistakes)
    \item Algebra becoming unmanageably complex
\end{itemize}

\textbf{Hard Abandon Triggers (Immediate switch):}
\begin{itemize}[leftmargin=*]
    \item After 6 minutes total time investment
    \item Realized fundamental conceptual error in approach
    \item Emotional frustration affecting clear thinking
\end{itemize}

\textbf{Recovery Strategy:}
\begin{itemize}[leftmargin=*]
    \item Mark answer (D) based on answer choice analysis (most complex = likely correct)
    \item Flag for review if time permits at end
    \item Move to next problem to maintain momentum
\end{itemize}
\end{competitionbox}

\begin{keyinsightbox}
\textbf{Key Insight}: The trapezoid decomposition method completely avoids finding the third side $PR$ of the triangle. By recognizing that:
\[[PP'Q'Q] + [QQ'R'R] = [PQR] + [PP'R'R]\]

we can solve for $[PQR]$ using only simple trapezoid area formulas and the Pythagorean theorem applied to height differences. This transforms a potentially complex Heron's formula calculation into straightforward arithmetic with radicals.
\end{keyinsightbox}

\begin{warningbox}
\textbf{Warning}: The most common error is incorrectly applying the Pythagorean theorem. Students often forget that the vertical distance in the trapezoid is the \textbf{difference} of the radii (not their sum), and the hypotenuse is the distance between centers (not the sum of radii, though they happen to equal in this case due to external tangency).

Example: For trapezoid $PP'Q'Q$:
\begin{itemize}
    \item $\times$ WRONG: $P'Q'^2 + 3^2 = 3^2$ (using sum of radii for height)
    \item $\checkmark$ CORRECT: $P'Q'^2 + 1^2 = 3^2$ (using difference of radii)
\end{itemize}
\end{warningbox}

\begin{tipbox}
\textbf{Pro Tip}: Answer choice analysis provides a quick confidence check. Choice (D) $\sqrt{6} - \sqrt{2}$ is the most complex expression, suggesting the problem requires multiple radical calculations. When you see your calculation naturally producing $2\sqrt{2}$ and $2\sqrt{6}$, this confirms you're on the right track. In AMC12, more difficult problems (P13-20) often have more complex answer choices.

Additionally, recognize that $\sqrt{6} - \sqrt{2} = \sqrt{2}(\sqrt{3} - 1)$, which could be useful if you need to verify the answer or if the problem asked for the answer in a different form.
\end{tipbox}

\section*{Final Answer}
\[\boxed{\textbf{(D) } \sqrt{6}-\sqrt{2}}\]

\section*{Confidence Assessment}
\begin{itemize}[leftmargin=*]
    \item \textbf{Confidence Level}: 95\% - Clean calculation with exact answer choice match
    \item \textbf{Verification Applied}: Level 4 Standard Verification - dimensional analysis, reasonableness check, algebraic verification all passed
    \item \textbf{Flag for Review}: No - high confidence, all verification checks passed
    \item \textbf{Time Spent}: Estimated 4.5 minutes (within target range)
\end{itemize}

\section*{Alternative Solution Summary}

\subsection*{Method 2: Coordinate Geometry with Shoelace Formula}

\textbf{Setup}: Place $P' = (0,0)$ on line $l$ (the $x$-axis)
\begin{itemize}
    \item $P = (0, 1)$
    \item $Q = (2\sqrt{2}, 2)$ (since $P'Q' = 2\sqrt{2}$)
    \item $R = (2\sqrt{2} + 2\sqrt{6}, 3)$ (since $Q'R' = 2\sqrt{6}$)
\end{itemize}

\textbf{Shoelace Formula}:
\begin{align*}
\text{Area} &= \frac{1}{2}|x_1(y_2-y_3) + x_2(y_3-y_1) + x_3(y_1-y_2)|\\
&= \frac{1}{2}|0(2-3) + 2\sqrt{2}(3-1) + (2\sqrt{2}+2\sqrt{6})(1-2)|\\
&= \frac{1}{2}|0 + 4\sqrt{2} - 2\sqrt{2} - 2\sqrt{6}|\\
&= \frac{1}{2}|2\sqrt{2} - 2\sqrt{6}|\\
&= \frac{1}{2} \cdot 2|\sqrt{2} - \sqrt{6}|\\
&= |\sqrt{2} - \sqrt{6}|\\
&= \sqrt{6} - \sqrt{2}
\end{align*}

This method produces the same answer with similar computational complexity.

\subsection*{Method 3: Heron's Formula (More Complex)}

\textbf{Step 1}: Find the third side $PR$
\begin{align*}
PR^2 &= (P'R')^2 + (r_R - r_P)^2\\
&= (2\sqrt{2} + 2\sqrt{6})^2 + (3-1)^2\\
&= (2\sqrt{2})^2 + 2 \cdot 2\sqrt{2} \cdot 2\sqrt{6} + (2\sqrt{6})^2 + 4\\
&= 8 + 8\sqrt{12} + 24 + 4\\
&= 36 + 8\sqrt{12}\\
&= 36 + 16\sqrt{3}
\end{align*}

Therefore: $PR = \sqrt{36 + 16\sqrt{3}} = 2\sqrt{9 + 4\sqrt{3}}$

\textbf{Step 2}: Apply Heron's formula with $a=3$, $b=5$, $c=2\sqrt{9+4\sqrt{3}}$

Semi-perimeter: $s = \frac{3 + 5 + 2\sqrt{9+4\sqrt{3}}}{2} = 4 + \sqrt{9+4\sqrt{3}}$

Area:
\begin{align*}
K &= \sqrt{s(s-a)(s-b)(s-c)}\\
&= \sqrt{(4+a)(4+a-3)(4+a-5)(4+a-2a)} \text{ where } a = \sqrt{9+4\sqrt{3}}\\
&= \sqrt{(4+a)(1+a)(a-1)(4-a)}\\
&= \sqrt{[(4+a)(4-a)][(1+a)(a-1)]}\\
&= \sqrt{(16-a^2)(a^2-1)}\\
&= \sqrt{(16-9-4\sqrt{3})(9+4\sqrt{3}-1)}\\
&= \sqrt{(7-4\sqrt{3})(8+4\sqrt{3})}\\
&= \sqrt{56 + 28\sqrt{3} - 32\sqrt{3} - 48}\\
&= \sqrt{8 - 4\sqrt{3}}
\end{align*}

\textbf{Step 3}: Simplify nested radical $\sqrt{8-4\sqrt{3}}$

Assume $\sqrt{8-4\sqrt{3}} = \sqrt{a} - \sqrt{b}$ where $a > b$

Then: $8 - 4\sqrt{3} = a + b - 2\sqrt{ab}$

Comparing: $a + b = 8$ and $2\sqrt{ab} = 4\sqrt{3}$, so $ab = 12$

Solving: $a = 6$, $b = 2$

Therefore: $\sqrt{8-4\sqrt{3}} = \sqrt{6} - \sqrt{2}$ $\checkmark$

This method is significantly more complex but serves as excellent verification.

\section*{Learning Outcomes}

\subsection*{Key Techniques Mastered}
\begin{enumerate}
    \item Pythagorean theorem in trapezoid configurations
    \item Area decomposition methods for triangles
    \item Radical simplification: $\sqrt{8} = 2\sqrt{2}$, $\sqrt{24} = 2\sqrt{6}$
    \item Nested radical simplification: $\sqrt{8-4\sqrt{3}} = \sqrt{6} - \sqrt{2}$
\end{enumerate}

\subsection*{Strategic Lessons}
\begin{enumerate}
    \item Multiple approaches exist - choose the cleanest calculation
    \item Trapezoid decomposition often cleaner than Heron's formula
    \item Answer choice complexity provides hints about solution path
    \item Verification is crucial for problems with multiple radical terms
\end{enumerate}

\subsection*{Competition Wisdom}
\begin{enumerate}
    \item P15 problems should take 4-6 minutes - budget accordingly
    \item Clean geometric diagrams prevent calculation errors
    \item Recognize when algebra is getting too messy - switch methods
    \item High confidence on mid-band problems justifies moving forward
\end{enumerate}

\end{document}